
\documentclass[11pt,oneside]{article}

\usepackage[T1]{fontenc}
\usepackage[utf8]{inputenc}
\usepackage{lmodern}
\usepackage{microtype}
\usepackage[a4paper,margin=29mm,headheight=14pt]{geometry}
\usepackage{amsmath,amssymb,amsthm,mathtools}
\usepackage{booktabs,longtable,array}
\usepackage{enumitem}
\usepackage{xcolor}
\usepackage{fancyhdr}
\usepackage{hyperref}
\usepackage[nameinlink,capitalise,noabbrev]{cleveref}

\definecolor{LinkBlue}{RGB}{20,63,120}
\definecolor{RuleGray}{RGB}{95,95,95}
\hypersetup{
  colorlinks=true,
  linkcolor=LinkBlue,
  citecolor=LinkBlue,
  urlcolor=LinkBlue,
  pdftitle={Operator-Theoretic Continuation on the Alaoglu--Erdos Conjecture},
  pdfauthor={OpenAI}
}

\pagestyle{fancy}
\fancyhf{}
\fancyhead[L]{Alaoglu--Erd\H{o}s: operator-theoretic continuation}
\fancyhead[R]{\thepage}
\renewcommand{\headrulewidth}{0.3pt}

\newtheorem{theorem}{Theorem}[section]
\newtheorem{proposition}[theorem]{Proposition}
\newtheorem{corollary}[theorem]{Corollary}
\newtheorem{lemma}[theorem]{Lemma}
\newtheorem{criterion}[theorem]{Closing criterion}
\theoremstyle{definition}
\newtheorem{definition}[theorem]{Definition}
\newtheorem{construction}[theorem]{Construction}
\newtheorem{question}[theorem]{Research question}
\theoremstyle{remark}
\newtheorem{remark}[theorem]{Remark}
\newtheorem{warning}[theorem]{Warning}

\newcommand{\N}{\mathbb N}
\newcommand{\Nzero}{\mathbb N_0}
\newcommand{\Z}{\mathbb Z}
\newcommand{\Q}{\mathbb Q}
\newcommand{\R}{\mathbb R}
\newcommand{\C}{\mathbb C}
\newcommand{\Qp}{\mathbb Q_{>0}}
\newcommand{\Alg}{\overline{\mathbb Q}}
\newcommand{\spn}{\operatorname{span}}
\newcommand{\supp}{\operatorname{supp}}
\newcommand{\Spec}{\operatorname{Spec}}
\newcommand{\Tr}{\operatorname{Tr}}
\newcommand{\rank}{\operatorname{rank}}
\newcommand{\diag}{\operatorname{diag}}
\newcommand{\ess}{\mathrm{ess}}
\newcommand{\floor}[1]{\left\lfloor #1\right\rfloor}
\newcommand{\ceil}[1]{\left\lceil #1\right\rceil}
\newcommand{\ee}{\mathrm e}
\newcommand{\ii}{\mathrm i}
\newcommand{\id}{\mathrm{id}}
\newcommand{\cA}{\mathcal A}
\newcommand{\cH}{\mathcal H}
\newcommand{\cF}{\mathcal F}
\newcommand{\cK}{\mathcal K}
\newcommand{\cT}{\mathcal T}
\newcommand{\cP}{\mathcal P}

\setlist[itemize]{leftmargin=1.6em,itemsep=0.3em,topsep=0.4em}
\setlist[enumerate]{leftmargin=1.8em,itemsep=0.35em,topsep=0.4em}
\setlength{\parindent}{0pt}
\setlength{\parskip}{0.58em}
\setlength{\emergencystretch}{3em}
\setcounter{tocdepth}{2}

\title{\textbf{Operator-Theoretic Continuation on the\\
Alaoglu--Erd\H{o}s Conjecture}\\[0.5em]
\large Local universality, arithmetic Hamiltonians, affine rigidity,\\
Jordan self-extensions, and Koopman realizations}
\author{Research continuation report}
\date{21 August 2026}

\begin{document}
\maketitle

\begin{abstract}
Let \(x\in\R\) and suppose that \(A=2^x\) and \(B=3^x\) are
integers.  The Alaoglu--Erd\H{o}s conjecture asserts that \(x\) is an
integer.  This document is a continuation of the supplied unified report,
not a replacement for it.  It deliberately avoids repeating that report's
historical survey, standard transcendence-theoretic reductions, extensive
catalogue of equivalent formulations, and previously explored analytic,
\(p\)-adic, matrix, and asymptotic routes.

The new organizing question is whether the scalar data can be forced to
extend to an operator structure carrying enough \emph{additive},
\emph{infinitesimal}, or \emph{finite-lattice} coherence to make a
nonintegral exponent impossible.  Several exact results are obtained.
Pairs of bilateral multiplication shifts are universal and therefore
spectrally blind to the arithmetic labels.  The two-mode Fock Hamiltonians
for \((2,3)\) and \((A,B)\) are exactly related by time rescaling, so heat
traces, local KMS data, and Weyl asymptotics supply no obstruction.  In
contrast, a global isometric intertwiner for the arithmetic Hamiltonian
\(H e_n=(\log n)e_n\) exists precisely for positive integral scaling
factors.  A finite prime-mode completion makes \(x\) an eigenvalue of an
integer valuation matrix and hence an integer; an automorphic completion
forces \(x=1\).  Adjoining even one affine \(ax+b\) coherence relation is
quantitatively rigid, including modulo compact operators.

Finite-dimensional functional calculus exhibits a complementary
dichotomy.  Every semisimple rational operator supported at the two
spectral points \(2\) and \(3\) is automatically blind: its \(x\)-th power
is rational whenever \(2^x\) and \(3^x\) are.  A single nontrivial Jordan
self-extension probes the derivative \(x2^{x-1}\) or \(x3^{x-1}\), and
rationality of that first jet forces \(x\in\Q\), hence \(x\in\Z\).
For the fractional part \(\alpha=\{x\}\), two rational-valued weighted
Koopman operators are constructed.  Their cocycles are bounded
coboundaries, explaining why spectra, Lyapunov growth, entropy, and
determinant-like invariants collapse.  The surviving datum is a projection
of trace \(\alpha\): an exact finite-dimensional tracial model or an exact
finite-Hankel realization would force \(\alpha\) rational, whereas ordinary
matrix approximation would not.

No unconditional proof of the conjecture is claimed.  The report instead
isolates a sharp local-to-global obstruction and ranks concrete closing
targets.  The strongest small targets are: propagation to one third
multiplicatively independent algebraic base; one rational Jordan
self-extension; one finite prime-mode completion; one exact finite-Hankel
realization; or one additive affine coherence relation.  Each target is
proved sufficient, and the seductive operator invariants that merely
repackage the original \(2\times2\) exponential array are explicitly
separated from those that add genuinely new arithmetic information.
\end{abstract}

\tableofcontents
\newpage

\section{Scope, conventions, and relation to the companion report}

The supplied unified report already develops the classical reduction to a
special case of the Four Exponentials Conjecture, the
Gelfond--Schneider and Six Exponentials thresholds, logarithmic and
valuation formulations, and a wide range of analytic and algebraic
approaches.  Repeating those discussions would obscure the purpose of a
continuation.  Here they are used only as black-box inputs when an
operator construction reaches one of their exact closing thresholds.

The main novelty is the distinction between three levels of structure.

\begin{enumerate}
\item \textbf{Local multiplicative structure.}  One retains only the two
      multiplicatively independent modes generated by \(2\) and \(3\).
      At this level, almost all natural operator invariants are universal
      or are changed only by a continuous rescaling of time.
\item \textbf{Finite-mode arithmetic structure.}  One asks whether the
      two modes extend to a finite lattice of prime modes, to a
      nonsemisimple rational functional calculus, or to a finite linear
      realization.  Any such exact extension has discrete algebraic
      invariants and is already strong enough to close the conjecture.
\item \textbf{Additive--multiplicative structure.}  One couples
      multiplication by \(n\) to translation by \(1\), as in the affine
      semigroup.  Integer labels then become rigid rather than merely
      spectral parameters.
\end{enumerate}

All Hilbert spaces are complex and separable.  For an unbounded operator,
an intertwining identity is initially interpreted on the algebraic span
of the indicated standard basis, which is a common invariant core.
The positive integers are \(\N=\{1,2,\ldots\}\), while
\(\Nzero=\{0,1,\ldots\}\).

\subsection{Minimal standing reduction}

Set
\[
        A=2^x\in\N,\qquad B=3^x\in\N.
\]
If \(x<0\), then \(0<2^x<1\), impossible for an integer; hence \(x\geq0\).
The case \(x=0\) is already integral, so the substantive case has \(x>0\).
Write
\[
        x=n+\alpha,\qquad n=\floor{x},\qquad 0\leq\alpha<1.
\]
Then
\begin{equation}\label{eq:r-s}
        r:=2^\alpha=\frac{A}{2^n}\in\Q,\qquad
        s:=3^\alpha=\frac{B}{3^n}\in\Q.
\end{equation}

\begin{lemma}[Rational-exponent collapse]\label{lem:rational-collapse}
If \(b\geq2\) is an integer, \(y\in\Q\), and \(b^y\in\Q\), then the
denominator of \(y\) divides every valuation of \(b\).  In particular,
if \(b\) is not a perfect power and \(b^y\in\Q\), then \(y\in\Z\).
Consequently, under \eqref{eq:r-s}, \(\alpha\in\Q\) implies
\(\alpha=0\).
\end{lemma}

\begin{proof}
Write \(y=p/q\) in lowest terms with \(q>0\), and write
\(b^y=u/v\) in lowest terms.  Taking \(q\)-th powers and applying each
prime valuation gives
\[
        q\bigl(v_\ell(u)-v_\ell(v)\bigr)=p\,v_\ell(b).
\]
Since \(\gcd(p,q)=1\), the integer \(q\) divides every \(v_\ell(b)\).
For \(b=2\), this gives \(q=1\).  Since \(0\leq\alpha<1\), one then has
\(\alpha=0\).
\end{proof}

Thus a counterexample would have \(x\) and \(\alpha\) irrational; by
Gelfond--Schneider it would in fact have \(x\) transcendental
\cite{Gelfond,Schneider,BakerBook}.  The point
of the operator constructions below is not to rediscover this reduction,
but to identify what additional exact structure would force the missing
rationality.

\subsection{A result ledger}

\begin{longtable}{>{\raggedright\arraybackslash}p{0.21\textwidth}
                  >{\raggedright\arraybackslash}p{0.49\textwidth}
                  >{\raggedright\arraybackslash}p{0.20\textwidth}}
\toprule
\textbf{Structure} & \textbf{Exact conclusion} & \textbf{Role}\\
\midrule
\endfirsthead
\toprule
\textbf{Structure} & \textbf{Exact conclusion} & \textbf{Role}\\
\midrule
\endhead
Bilateral shifts for two independent labels &
All such pairs are unitarily equivalent with countable multiplicity. &
No-go: pure spectrum is blind.\\
Two-mode Fock Hamiltonian &
\(H_{A,B}=xH_{2,3}\); heat traces and local KMS data are time-rescaled. &
No-go: thermodynamics is blind.\\
Arithmetic Hamiltonian on all \(\N\) &
An isometry \(T\) with \(HT=xTH\) exists exactly for
\(x\in\N\). &
Global closing criterion.\\
Finite invariant prime lattice &
\(x\) is an algebraic-integer eigenvalue of a valuation matrix, hence
\(x\in\Z\). &
Finite-mode closing criterion.\\
Unimodular prime-lattice completion &
The same eigenvalue is an algebraic unit; hence \(x=1\). &
Automorphic closing criterion.\\
Affine \(ax+b\) coherence &
Preserving \(V_2S=S^2V_2\), even modulo compacts in the stated model,
forces \(A=2\), hence \(x=1\). &
Robust additive rigidity.\\
Semisimple rational functional calculus &
Every operator with minimal polynomial dividing
\((t-2)(t-3)\) remains rational under \(t\mapsto t^x\). &
No-go: two values suffice.\\
Jordan self-extension &
Rationality of \((2I+N)^x\), \(N^2=0\neq N\), forces \(x\in\Q\). &
Infinitesimal closing criterion.\\
Sturmian Koopman realization &
The rational weights are bounded coboundaries; one projection has trace
\(\alpha\). &
Separates blind invariants from trace.\\
Exact finite tracial or Hankel model &
The trace/frequency \(\alpha\) is rational, so \(\alpha=0\). &
Finite-realization closing criterion.\\
One third independent algebraic base &
Six Exponentials forces \(x\in\Q\), hence \(x\in\Z\). &
Smallest high-value propagation target.\\
\bottomrule
\end{longtable}

\section{The maximal two-mode algebraic sector}

Before introducing operator models, it is useful to sharpen exactly how
far the hypotheses propagate by known transcendence theory.  This result
is the correct boundary condition for every later local-to-global
construction.

\begin{theorem}[Three-mode barrier]\label{thm:three-mode}
Assume \(x\notin\Q\), \(2^x,3^x\in\Alg\), and let \(q\in\Qp\).
Then
\[
        q^x\in\Alg
        \quad\Longleftrightarrow\quad
        q=2^m3^n\ \text{for some }m,n\in\Z.
\]
In particular, if a nonintegral solution of the Alaoglu--Erd\H{o}s
hypotheses exists, then
\[
 \{k\in\N:k^x\in\Alg\}
 =
 \{2^m3^n:m,n\in\Nzero\}.
\]
Every positive integer containing a prime other than \(2\) or \(3\) has
transcendental \(x\)-th power.
\end{theorem}

\begin{proof}
The reverse implication is immediate:
\((2^m3^n)^x=(2^x)^m(3^x)^n\) is algebraic.

For the forward implication, suppose that \(q\) is multiplicatively
independent of \(2\) and \(3\).  Then
\(\log2,\log3,\log q\) are linearly independent over \(\Q\), while
\(1,x\) are linearly independent over \(\Q\).  The Six Exponentials
Theorem \cite[Chap.~3]{WaldschmidtBook} says that at least one entry of
\[
\begin{pmatrix}
 \exp(\log2) & \exp(\log3) & \exp(\log q)\\
 \exp(x\log2) & \exp(x\log3) & \exp(x\log q)
\end{pmatrix}
=
\begin{pmatrix}
 2&3&q\\ 2^x&3^x&q^x
\end{pmatrix}
\]
is transcendental, contradicting the assumed algebraicity of all six
entries.  Therefore \(q\) is multiplicatively dependent with \(2,3\):
there are integers \(a,b,c\), with \(c\neq0\), such that
\(q^c2^a3^b=1\).  Prime valuations show that \(q\) has no prime factor
other than \(2\) or \(3\); since \(q\in\Qp\), it follows that
\(q=2^m3^n\) with \(m,n\in\Z\).
\end{proof}

\begin{remark}
The theorem says more than ``a third prime would help.''  Under a
hypothetical counterexample, the rank-two group
\(\langle2,3\rangle\) is the \emph{entire} rational domain on which the
\(x\)-power map remains algebraic.  An operator proof must therefore
force leakage out of this maximal sector, or extract an infinitesimal or
finite-dimensional invariant that cannot live inside it.
\end{remark}

\begin{corollary}[One-vector closing target]\label{cor:one-vector}
Under the Alaoglu--Erd\H{o}s hypotheses, it suffices to prove
\(q^x\in\Alg\) for one rational \(q>0\) having a prime factor other than
\(2\) or \(3\).  In particular, any one of \(5^x,7^x,10^x,\ldots\)
being algebraic closes the conjecture.
\end{corollary}

\section{Bilateral multiplication shifts: a universality theorem}

Let \(G=\Qp\) with multiplication, and let
\(\cH_G=\ell^2(G)\) with standard basis \((e_q)_{q\in G}\).  For
\(a\in G\), define a unitary
\[
        U_a e_q=e_{aq}.
\]
These are the most literal bilateral operator representatives of
multiplication by rational labels; the orbit decomposition below is the
regular-representation viewpoint standard in semigroup operator theory
\cite{Nica,LacaRaeburn}.

\begin{theorem}[Universality of two independent multiplication shifts]
\label{thm:bilateral-universal}
If \(a,b\in\Qp\) are multiplicatively independent, then the commuting
pair \((U_a,U_b)\) is unitarily equivalent to a countably infinite
multiple of the left regular representation of \(\Z^2\).  Consequently,
for any two multiplicatively independent pairs \((a,b)\) and \((c,d)\),
the pairs \((U_a,U_b)\) and \((U_c,U_d)\) are unitarily equivalent.
\end{theorem}

\begin{proof}
Let \(\Lambda=\langle a,b\rangle=\{a^mb^n:m,n\in\Z\}\).
Multiplicative independence identifies \(\Lambda\) with \(\Z^2\).
Choose a transversal \(C\) for the cosets \(G/\Lambda\).  Then
\[
        \ell^2(G)=\bigoplus_{c\in C}\ell^2(c\Lambda),
\]
and on each summand the unitary
\[
        e_{ca^mb^n}\longmapsto e_{(m,n)}
\]
conjugates \(U_a,U_b\) to the two coordinate bilateral shifts on
\(\ell^2(\Z^2)\).  The quotient \(G/\Lambda\) is countably infinite:
it is countable, and primes outside the finite prime support of \(ab\)
give infinitely many distinct cosets.  Thus the multiplicity is
countably infinite and does not depend on the labels.
\end{proof}

For \(x>0\), the integers \(A=2^x\) and \(B=3^x\) are
multiplicatively independent.  Indeed,
\(A^mB^n=1\) implies \((2^m3^n)^x=1\), hence \(m=n=0\).
Therefore \cref{thm:bilateral-universal} applies equally to
\((2,3)\) and \((A,B)\).

\begin{corollary}[Spectral blindness]\label{cor:spectral-blindness}
The joint spectrum, joint spectral type, von Neumann algebra generated by
the pair, and spectral multiplicity of the bilateral pair do not
distinguish \((2,3)\) from \((2^x,3^x)\).  Any proposed proof based only
on the unitary-equivalence class of the two multiplication shifts cannot
see whether \(x\) is integral.
\end{corollary}

\begin{remark}
Under Fourier transform, one copy of the regular representation becomes
multiplication by the two coordinate functions on \(L^2(\mathbb T^2)\).
Thus the joint spectrum is always \(\mathbb T^2\), with the same Haar spectral
measure and the same countable multiplicity.  The failure is not a lack
of a sufficiently clever spectral invariant; the entire pair is
universal.
\end{remark}

\subsection{The logarithmic Hamiltonian restores scale locally}

Define the self-adjoint arithmetic Hamiltonian
\[
        H_G e_q=(\log q)e_q.
\]
On its natural core,
\begin{equation}\label{eq:comm-H-U}
        [H_G,U_a]=(\log a)U_a.
\end{equation}
The commutator remembers the logarithmic label that pure shift spectrum
forgets.  Nevertheless, on the two generated orbits it still gives no
obstruction.

Let
\[
 \cH_{2,3}=\overline{\spn}\{e_{2^m3^n}:m,n\in\Z\},\qquad
 \cH_{A,B}=\overline{\spn}\{e_{A^mB^n}:m,n\in\Z\}.
\]
The map
\begin{equation}\label{eq:local-W}
        W e_{2^m3^n}=e_{A^mB^n}
\end{equation}
is a unitary \(\cH_{2,3}\to\cH_{A,B}\), and
\begin{equation}\label{eq:local-H-cov}
        H_GW=xWH_G.
\end{equation}
Thus even the Hamiltonian-covariant equivalence is automatic on the
maximal rank-two algebraic sector.  The arithmetic content lies in
extending \eqref{eq:local-W} beyond that sector while preserving a
discrete basis or lattice.

\section{Two-mode Fock Hamiltonians and thermodynamic no-go results}

Let \(\cF_2=\ell^2(\Nzero^2)\), with number operators
\[
        N_1e_{m,n}=m e_{m,n},\qquad
        N_2e_{m,n}=n e_{m,n}.
\]
For multiplicatively independent \(a,b>1\), set
\[
        H_{a,b}=(\log a)N_1+(\log b)N_2.
\]
Its eigenvalue at \(e_{m,n}\) is
\(\log(a^mb^n)\), so the spectrum is simple.

\begin{proposition}[Exact time homothety]\label{prop:fock-homothety}
Under \(A=2^x\) and \(B=3^x\),
\[
        H_{A,B}=xH_{2,3}
\]
as self-adjoint operators on the same Fock space.
\end{proposition}

This equality is stronger than an asymptotic similarity.  Every spectral
quantity derived functorially from \(H\) is merely reparametrized.

For \(\beta>0\), the heat trace is
\begin{equation}\label{eq:partition}
 Z_{a,b}(\beta)
 :=\Tr(\ee^{-\beta H_{a,b}})
 =\sum_{m,n\geq0}a^{-\beta m}b^{-\beta n}
 =\frac{1}{(1-a^{-\beta})(1-b^{-\beta})}.
\end{equation}
Consequently,
\[
        Z_{A,B}(\beta)=Z_{2,3}(x\beta).
\]
As \(\beta\downarrow0\),
\begin{equation}\label{eq:heat-asymp}
        Z_{a,b}(\beta)
        =\frac{1}{\beta^2\log a\,\log b}
         +O(\beta^{-1}),
\end{equation}
and the eigenvalue-counting function satisfies the elementary lattice
estimate
\begin{equation}\label{eq:weyl}
 N_{a,b}(E)
 :=\#\{(m,n)\in\Nzero^2:m\log a+n\log b\leq E\}
 =\frac{E^2}{2\log a\,\log b}+O(E).
\end{equation}
Both leading terms simply acquire the continuous factor \(x^{-2}\) when
\((a,b)=(2,3)\) is replaced by \((A,B)\).

Let \(L_1,L_2\) be the coordinate unilateral shifts on \(\cF_2\).
The Toeplitz algebra \(\cT^{(2)}=C^*(L_1,L_2)\) admits the dynamics
\[
        \sigma_t^{a,b}(L_1)=a^{\ii t}L_1,\qquad
        \sigma_t^{a,b}(L_2)=b^{\ii t}L_2.
\]
Then
\[
        \sigma_t^{A,B}=\sigma_{xt}^{2,3}.
\]
Thus the underlying \(C^*\)-algebra, its \(K\)-theory, and the local
two-mode representation are identical; the dynamics differs only by a
time change.  Gibbs and KMS data parameterized by inverse temperature
are transported by \(\beta\mapsto x\beta\); see
\cite{BratteliRobinson,BostConnes} for the general equilibrium formalism.

\begin{warning}[Why a critical temperature cannot quantize \(x\)]
A critical inverse temperature is meaningful only after fixing an
external normalization of time or energy.  In the isolated two-mode
system there is no such normalization: multiplying the Hamiltonian by
an arbitrary positive real number is an exact time reparametrization.
Therefore a pole, heat coefficient, entropy density, or KMS transition
inside this sector cannot force \(x\in\Z\).
\end{warning}

\section{The full prime Fock space and the local-to-global gap}

Let \(\cP\) be the set of primes and
\[
        \Nzero^{(\cP)}
        =\{v=(v_p)_{p\in\cP}:v_p\in\Nzero,\ 
          v_p=0\text{ for all but finitely many }p\}.
\]
Unique factorization identifies \(\N\) with this monoid through the
valuation vector
\[
        \nu(n)=(v_p(n))_{p\in\cP}.
\]
The full prime Fock space is
\[
        \cF_{\cP}=\ell^2(\Nzero^{(\cP)})\cong\ell^2(\N),
\]
and its Hamiltonian is
\[
        H_{\cP}e_v
        =\left(\sum_{p\in\cP}v_p\log p\right)e_v.
\]
Under the identification with \(\ell^2(\N)\), this is simply
\(H e_n=(\log n)e_n\).

Let \(a=\nu(A)\) and \(b=\nu(B)\).  Since \(A,B\) are multiplicatively
independent, the map
\[
        J:\cF_2\longrightarrow\cF_{\cP},\qquad
        J e_{m,n}=e_{ma+nb}
\]
is an isometry.  Moreover,
\begin{equation}\label{eq:fock-embedding}
        H_{\cP}J=xJH_{2,3}.
\end{equation}
Thus the hypotheses themselves give an exact Hamiltonian-covariant
embedding of the two-mode sector into the arithmetic spectrum.  This is
not a partial proof; it is the canonical operator form of the two scalar
equalities.

The missing step is now precise:

\begin{quote}
\emph{Can the two-column valuation map
\(e_2\mapsto\nu(A)\), \(e_3\mapsto\nu(B)\) be forced to extend to a
finite or global integral lattice map that rescales the logarithmic
energy functional by \(x\)?}
\end{quote}

Abstract Hilbert-space extension theorems are irrelevant here.  The
isometry \(J\) can always be extended abstractly after adding arbitrary
orthogonal complements.  What matters is an extension that remains
\emph{monomial on the arithmetic basis}, \emph{integral on valuation
lattices}, and \emph{covariant for \(H_{\cP}\)}.  The next two sections
show that either a global extension or a finite prime-mode extension
would close the conjecture.


\section{Global arithmetic-Hamiltonian intertwiners}

On \(\ell^2(\N)\), let
\[
        He_n=(\log n)e_n,
\]
with domain
\[
        \mathcal D(H)
        =\left\{\xi:\sum_{n\geq1}(\log n)^2|\xi_n|^2<\infty\right\}.
\]
The eigenvalues are simple.  This converts a Hamiltonian covariance
identity into an exact arithmetic condition on every basis label.

\begin{theorem}[Global isometric covariance]\label{thm:global-isometry}
For \(c>0\), the following are equivalent.
\begin{enumerate}
\item There is an isometry \(T:\ell^2(\N)\to\ell^2(\N)\) such that
      \(HT=cTH\) on the finite-support core.
\item \(n^c\in\N\) for every \(n\in\N\).
\item \(c\in\N\).
\end{enumerate}
For \(c=k\in\N\), a canonical intertwiner is
\[
        T_k e_n=e_{n^k}.
\]
\end{theorem}

\begin{proof}
Suppose \(HT=cTH\).  For each \(n\),
\[
        H(Te_n)=c(\log n)Te_n.
\]
Since \(T\) is an isometry, \(Te_n\neq0\).  The spectrum of \(H\) is
simple and consists of the numbers \(\log m\), \(m\in\N\), so there is
an \(m\in\N\) with
\(\log m=c\log n\); hence \(m=n^c\in\N\).  This proves (1)\(\Rightarrow\)(2).

To prove (2)\(\Rightarrow\)(3), suppose that \(c\notin\N\), and choose
\(r=\floor{c}+1\).  The \(r\)-th forward difference
\[
        \Delta^r n^c
        =\sum_{j=0}^{r}(-1)^{r-j}\binom rj(n+j)^c
\]
is an integer.  Repeated use of the fundamental theorem of calculus
gives the exact integral formula
\[
 \Delta^r n^c
 =c(c-1)\cdots(c-r+1)
  \int_{[0,1]^r}
  (n+t_1+\cdots+t_r)^{c-r}\,dt_1\cdots dt_r.
\]
The coefficient is nonzero because \(c\) is not an integer, and the
integrand has fixed sign.  Thus the difference is nonzero for every
\(n\), while \(c-r<0\) implies \(\Delta^r n^c\to0\).  For large \(n\)
this is a nonzero integer of absolute value less than \(1\), a
contradiction.  Hence \(c\in\N\).

Finally, if \(c=k\in\N\), then \(n\mapsto n^k\) is injective, so \(T_k\)
is an isometry, and
\[
        HT_ke_n=k(\log n)e_{n^k}=kT_kHe_n.
\]
\end{proof}

\begin{remark}[The endpoint \(c=0\)]
The scalar map \(n\mapsto n^0=1\) collapses all basis vectors and cannot
be implemented by an isometry.  It is represented by a rank-one
collapse rather than an endomorphism.  This is why
\cref{thm:global-isometry} classifies positive integers; the original
conjecture separately includes the already trivial solution \(x=0\).
\end{remark}

\subsection{Partial intertwiners and exact support}

For arbitrary \(c>0\), define
\[
        I_c=\{n\in\N:n^c\in\N\}.
\]
There is a canonical isometry
\[
        T_c:\ell^2(I_c)\longrightarrow\ell^2(\N),\qquad
        T_c e_n=e_{n^c},
\]
and it satisfies \(HT_c=cT_cH\) on the natural source core.
Under the Alaoglu--Erd\H{o}s hypotheses,
\(\{2^m3^n:m,n\in\Nzero\}\subseteq I_x\), so an infinite-rank
Hamiltonian intertwiner already exists.

Combining this observation with \cref{thm:three-mode} gives a complete
classification under a hypothetical counterexample.

\begin{corollary}[Maximal support of the partial intertwiner]
\label{cor:max-partial-support}
If \(x\) is a nonintegral solution of the Alaoglu--Erd\H{o}s hypotheses,
then
\[
        I_x=\{2^m3^n:m,n\in\Nzero\}.
\]
More strongly, replacing ``integer'' by ``algebraic'' in the definition
of \(I_x\) gives the same set.
\end{corollary}

This identifies an especially small operator target: it is enough to
make the covariant partial intertwiner nonzero at \emph{one} basis vector
\(e_q\) with \(q\) not \(3\)-smooth and with algebraic target label.
No global extension is needed; \cref{cor:one-vector} then closes the
conjecture.

\section{Finite prime-mode closure and valuation matrices}

The global criterion in \cref{thm:global-isometry} is much stronger than
the original hypotheses.  A finite-dimensional lattice completion is
already enough.

Let \(P=\{p_1,\ldots,p_d\}\) be a finite set of primes containing \(2\).
Write
\[
        \Gamma_P
        =\left\{\prod_{j=1}^d p_j^{v_j}:v_j\in\Z\right\}
        \cong\Z^d.
\]

\begin{theorem}[Finite prime closure]\label{thm:finite-prime-closure}
Let \(c\in\R\setminus\{0\}\).  Suppose that
\begin{enumerate}
\item \(2\in P\);
\item \(p_j^c\in\Gamma_P\) for every \(p_j\in P\); and
\item \(2^c\) is algebraic (in particular, rational or integral).
\end{enumerate}
Then \(c\in\Z\).  If \(c\geq0\), then \(c\in\Nzero\).
\end{theorem}

\begin{proof}
For each \(j\), write uniquely
\[
        p_j^c=\prod_{i=1}^d p_i^{m_{ij}},
        \qquad m_{ij}\in\Z,
\]
and let \(M=(m_{ij})\in M_d(\Z)\).  Set
\[
        \ell=(\log p_1,\ldots,\log p_d)^{\mathsf T}.
\]
Taking real logarithms column by column gives
\[
        M^{\mathsf T}\ell=c\ell.
\]
Thus \(c\) is an eigenvalue of an integer matrix and therefore an
algebraic integer.  If \(c\) were algebraic irrational,
Gelfond--Schneider would make \(2^c\) transcendental, contrary to
assumption.  Hence \(c\in\Q\).  A rational algebraic integer is an
integer.
\end{proof}

\begin{corollary}[Automorphic closure]\label{cor:automorphic-closure}
In the setting of \cref{thm:finite-prime-closure}, suppose additionally
that the power map
\[
        \Gamma_P\longrightarrow\Gamma_P,\qquad q\longmapsto q^c,
\]
is onto.  Then \(M\in\mathrm{GL}_d(\Z)\), so \(c\) is an algebraic unit.
If \(c>0\), then \(c=1\).
\end{corollary}

\begin{proof}
Surjectivity makes the valuation matrix unimodular.  Every eigenvalue of
a unimodular integer matrix is an algebraic unit.  By
\cref{thm:finite-prime-closure}, \(c\) is a rational integer; the only
positive rational integer that is an algebraic unit is \(1\).
\end{proof}

\subsection{Operator form of the valuation theorem}

On \(\ell^2(\Z^d)\), define
\[
        H_\ell e_v=\langle\ell,v\rangle e_v.
\]
The power map on \(\Gamma_P\) is injective for \(c\neq0\), hence the
integer matrix \(M\) is injective on \(\Z^d\).  Therefore
\[
        T_Me_v=e_{Mv}
\]
is an isometry, and
\[
        H_\ell T_M=cT_MH_\ell
        \quad\Longleftrightarrow\quad
        M^{\mathsf T}\ell=c\ell.
\]
Thus the theorem may be read as follows: a finite arithmetic lattice
model of Hamiltonian time rescaling makes the scaling factor an
algebraic integer.  The scalar algebraicity \(2^c\in\Alg\) then reduces
it to an ordinary integer.

If all \(p_j^c\) are integers supported on \(P\), then \(M\) has
nonnegative entries and the same statement can be realized on the
positive Fock cone \(\ell^2(\Nzero^d)\).  If the matrix is irreducible,
the positive logarithm vector exhibits \(c\) as its Perron--Frobenius
eigenvalue.  Perron--Frobenius theory does not itself prove integrality;
the decisive extra input remains Gelfond--Schneider.

\begin{criterion}[Finite propagation closure]\label{crit:finite-propagation}
Starting from \(P_0=\{2,3\}\), suppose one can prove that \(p^x\) is
rational for every prime introduced as a divisor of a previously known
\(q^x\), and that the resulting prime-propagation process terminates in
a finite set \(P\).  Then \(P\) is invariant under the \(x\)-power map,
so \cref{thm:finite-prime-closure} proves \(x\in\Z\).
\end{criterion}

A counterexample therefore cannot admit any finite exact
prime-substitution completion.  This gives a concrete diagnostic for
operator correspondences, graph \(C^*\)-algebras, fusion matrices, and
substitution systems: if their finite incidence matrix really encodes
the logarithmic scaling, the conjecture is already over.

\section{\(C^*\)-dynamical interpretation: local covariance versus global arithmetic}

The preceding Hilbert-space constructions have a direct
\(C^*\)-dynamical reading.  It is useful to state it without relying on
a particular presentation of the Bost--Connes algebra.

Let \(\mathfrak T_{\mathrm{mult}}\) be the concrete Toeplitz algebra on
\(\ell^2(\N)\) generated by the multiplication isometries
\[
        V_m e_n=e_{mn},\qquad m,n\in\N.
\]
They satisfy \(V_mV_n=V_{mn}\).  The Hamiltonian \(H e_n=(\log n)e_n\)
implements the time evolution
\begin{equation}\label{eq:time-evolution}
        \sigma_t(V_m)=m^{\ii t}V_m.
\end{equation}

On the algebraic, nonselfadjoint semigroup generated by
\(V_2,V_3\), the assignment
\[
        V_2\longmapsto V_A,\qquad
        V_3\longmapsto V_B
\]
respects products and satisfies
\[
        \sigma_t\!\left(\Phi(V_j)\right)
        =\Phi\!\left(\sigma_{xt}(V_j)\right),
        \qquad j=2,3.
\]
At the \(C^*\)-level one should use the two-coordinate Fock compression
from \cref{prop:fock-homothety}: adjoints and range projections in the
full concrete representation can additionally record gcd/lcm data.
That extra data is precisely valuation-lattice data; it does not follow
from abstract rank-two dynamics alone.  In either formulation, local
time covariance merely rescales the two known energies and does not
quantize \(x\).

A global monomial endomorphism would have to send each \(V_n\) to
\(V_{n^x}\).  Such labels exist for all \(n\) exactly when
\(n^x\in\N\) for all \(n\), and \cref{thm:global-isometry} then forces
\(x\in\N\).  Likewise, a finite prime-generated endomorphism gives the
integer matrix of \cref{thm:finite-prime-closure}.

\begin{remark}[Bost--Connes and related systems]
In Bost--Connes-type systems \cite{BostConnes}, the multiplicative isometries carry the
same logarithmic time evolution, while additional additive or torsion
data encode genuinely arithmetic relations.  Restricting to the two
multiplicative modes loses precisely that extra coherence and leaves
only a time-rescaled rank-two system.  A useful operator strategy must
therefore exploit the additive diagonal, transfer operators, or an
integral lattice completion; KMS data from the isolated modes cannot do
the job.
\end{remark}

\section{Affine \(ax+b\) rigidity}

The cleanest demonstration that additive structure restores integer
labels uses the standard affine representation familiar from
\(ax+b\) semigroup \(C^*\)-algebras \cite{CuntzAxB,BrownloweEtAl} on
\(\ell^2(\Nzero)\).  Define
\[
        Se_k=e_{k+1},\qquad
        V_m e_k=e_{mk}.
\]
Then
\begin{equation}\label{eq:axb}
        V_mS=S^mV_m.
\end{equation}
Unlike the purely multiplicative relations, this relation records the
integer \(m\) as an actual translation length.

\begin{theorem}[Quantitative affine rigidity]\label{thm:affine-rigidity}
For integers \(a,b,d\geq1\), set
\[
        D_{a,b,d}=V_aS^d-S^{bd}V_a.
\]
Then
\[
        D_{a,b,d}=0\quad\Longleftrightarrow\quad a=b.
\]
If \(a\neq b\), then \(D_{a,b,d}\) is not compact and
\[
        \|D_{a,b,d}\|_{\mathrm{ess}}\geq\sqrt2.
\]
In particular, the affine relation determines the multiplicative label
even in the Calkin algebra and is stable under norm errors smaller than
\(\sqrt2\).
\end{theorem}

\begin{proof}
For every \(k\geq0\),
\[
 D_{a,b,d}e_k
 =e_{a(k+d)}-e_{ak+bd}
 =e_{ak+ad}-e_{ak+bd}.
\]
If \(a=b\), this is zero.  If \(a\neq b\), the two displayed basis
vectors are distinct, so
\[
        \|D_{a,b,d}e_k\|=\sqrt2
\]
for every \(k\).  Since \(e_k\to0\) weakly and every compact operator
maps this sequence to a norm-null sequence, for every compact \(K\),
\[
        \liminf_{k\to\infty}\|(D_{a,b,d}-K)e_k\|\geq\sqrt2.
\]
Taking the infimum over \(K\) proves the essential-norm bound and
noncompactness.
\end{proof}

\begin{corollary}[One affine coherence closes the conjecture]
\label{cor:one-affine}
Suppose a proposed operator extension sends
\[
        S\longmapsto S^d,\qquad V_2\longmapsto V_A
\]
for some \(d\geq1\), and preserves the single relation
\(V_2S=S^2V_2\), either exactly or modulo compact operators.  Then
\(A=2\), hence \(x=1\).
\end{corollary}

\begin{proof}
The image of the relation requires
\(V_AS^d-S^{2d}V_A\) to vanish, or at least to be compact.
Apply \cref{thm:affine-rigidity} with \(a=A\) and \(b=2\).
\end{proof}

\begin{remark}
This criterion is intentionally much stronger than the original scalar
hypotheses, but it is structurally minimal: one generator \(S\), one
multiplication mode \(V_2\), and one coherence square suffice.  It also
explains why a proof cannot emerge from the multiplicative semigroup
alone.  Translation gives an external integral ruler against which the
label \(A\) can be measured.
\end{remark}

\subsection{A categorical shadow}

Let \(E_n\) denote an \(n\)-dimensional object in a semisimple tensor
category with direct sums.  The elementary coherence
\[
        E_2\otimes E_2\cong E_2\oplus E_2
\]
expresses \(2\cdot2=2+2\).  If a dimension-scaling construction sends
\(E_2\) to an object of dimension \(A=2^x\) and preserves this one
isomorphism, then
\[
        A^2=2A.
\]
Since \(A>0\), one gets \(A=2\) and \(x=1\).  The statement is simple,
but its lesson is exact: a single additive--multiplicative coherence is
vastly more rigid than all multiplicative powers of the two generators.


\section{Finite-dimensional functional calculus: semisimple blindness}

A direct matrix formulation replaces the two scalars by
\(\diag(2,3)\) and observes that
\(\diag(2,3)^x=\diag(A,B)\).  By itself this is exactly the original
problem.  More revealing is the behavior of all rational operators
whose spectrum is supported at \(2\) and \(3\).

Let \(f_x(t)=t^x\) on the positive real axis, extended holomorphically
to the slit plane; standard matrix-function facts used below may be
found in \cite{Higham} \(\C\setminus(-\infty,0]\).

\begin{theorem}[Semisimple two-point calculus is automatic]
\label{thm:semisimple-blind}
Assume \(A=2^x\in\Q\) and \(B=3^x\in\Q\).  If
\(T\in M_d(\Q)\) satisfies
\[
        (T-2I)(T-3I)=0,
\]
then
\begin{equation}\label{eq:semisimple-formula}
        T^x=(B-A)T+(3A-2B)I\in M_d(\Q).
\end{equation}
If \(T,A,B\) have integer entries, then \(T^x\) has integer entries.
\end{theorem}

\begin{proof}
The minimal polynomial is square-free, so \(T\) is semisimple with
spectral projections
\[
        P_2=3I-T,\qquad P_3=T-2I.
\]
Since \(P_2+P_3=I\),
\[
        T^x=AP_2+BP_3
        =A(3I-T)+B(T-2I),
\]
which is \eqref{eq:semisimple-formula}.
\end{proof}

The theorem covers every rationally diagonalizable coupling of the two
modes, not merely the diagonal matrix.  For example,
\[
        T_c=\begin{pmatrix}2&c\\0&3\end{pmatrix}
        \quad(c\in\Z)
\]
is diagonalizable and
\[
        T_c^x=
        \begin{pmatrix}
          A&c(B-A)\\0&B
        \end{pmatrix}.
\]
Thus adding an off-diagonal interaction between the distinct eigenvalues
does not add arithmetic content.  It merely evaluates the divided
difference
\[
        f_x[2,3]=\frac{3^x-2^x}{3-2}=B-A,
\]
which the original two scalar values already determine.

\begin{remark}[Module-theoretic interpretation]
Finite-dimensional \(\Q[t]\)-modules supported at the distinct maximal
ideals \((t-2)\) and \((t-3)\) split because the ideals are coprime.
Functional calculus on a semisimple module depends only on the two
values \(f_x(2)\) and \(f_x(3)\).  No extension class remains on which a
derivative could be detected.
\end{remark}

\section{Jordan self-extensions and the first-jet obstruction}

A repeated spectral point changes the situation.  Let \(N\neq0\) be a
rational matrix with \(N^2=0\).  The binomial series terminates:
\begin{equation}\label{eq:jordan-power}
        (bI+N)^x=b^xI+xb^{x-1}N.
\end{equation}

\begin{theorem}[One rational Jordan lift closes the conjecture]
\label{thm:jordan-closes}
Assume \(b\geq2\), \(b^x\in\Q\), and \(N\in M_d(\Q)\) satisfies
\(N^2=0\neq N\).  If \((bI+N)^x\in M_d(\Q)\), then \(x\in\Q\).
Consequently, under the Alaoglu--Erd\H{o}s hypotheses, rationality of
either \((2I+N)^x\) or \((3I+N)^x\) forces \(x\in\Z\).
\end{theorem}

\begin{proof}
Choose an entry \(N_{ij}\neq0\).  From
\eqref{eq:jordan-power}, rationality of the corresponding entry of
\((bI+N)^x\) gives \(xb^{x-1}N_{ij}\in\Q\), hence
\(xb^{x-1}\in\Q\).  Since \(b^x\in\Q^\times\),
\[
        x=\frac{b}{b^x}\,xb^{x-1}\in\Q.
\]
Under the Alaoglu--Erd\H{o}s hypotheses,
\cref{lem:rational-collapse} then gives \(x\in\Z\).
\end{proof}

\begin{criterion}[Minimal infinitesimal bridge]\label{crit:first-jet}
It is enough to prove that the scalar rationality \(2^x\in\Q\) extends
to the first-order dual-number identity
\[
        (2+\varepsilon)^x
        =2^x+x2^{x-1}\varepsilon
        \in\Q[\varepsilon]/(\varepsilon^2).
\]
Equivalently, it is enough to rationalize one nonzero tangent direction
at the spectral point \(2\) (or at \(3\)).
\end{criterion}

For a Jordan block of length \(r\),
\[
 (bI+N)^x
 =\sum_{j=0}^{r-1}\binom{x}{j}b^{x-j}N^j.
\]
The first superdiagonal already reveals \(xb^{x-1}\); higher jets add no
logical strength for rationality, although their divisibility may be
useful in attempts to derive the first jet from an integral model.

\subsection{The Loewner matrix diagnosis}

Let \(D=\diag(2,3)\).  The Fr\'echet derivative of \(f_x\) at \(D\) is
entrywise multiplication by the Loewner matrix:
\[
        \bigl(Df_x(D)[E]\bigr)_{ij}
        =f_x^{[1]}(\lambda_i,\lambda_j)E_{ij},
        \qquad(\lambda_1,\lambda_2)=(2,3),
\]
where
\[
 f_x^{[1]}(\lambda,\mu)=
 \begin{cases}
 \dfrac{\lambda^x-\mu^x}{\lambda-\mu},&\lambda\neq\mu,\\[1.2ex]
 x\lambda^{x-1},&\lambda=\mu.
 \end{cases}
\]
Under the scalar hypotheses, the off-diagonal coefficients are
\[
        f_x^{[1]}(2,3)=f_x^{[1]}(3,2)=B-A\in\Z,
\]
whereas the diagonal coefficients are
\[
        x2^{x-1},\qquad x3^{x-1}.
\]
This is a precise separation:

\begin{quote}
Coupling the \(2\)-mode to the \(3\)-mode probes only a secant, already
fixed by \(A\) and \(B\).  A self-coupling at either mode probes a
derivative and immediately forces \(x\) rational if that derivative is
rational.
\end{quote}

In extension-theoretic language, extensions between the two distinct
simple spectral modules split, while a self-extension is nontrivial and
carries the missing first-jet datum.

\subsection{A useful negative theorem}

\begin{proposition}[No semisimple matrix amplification can help]
\label{prop:no-semisimple-amplification}
Fix any finite family of rational matrices \(T_j\) whose minimal
polynomials are square-free and whose spectra are subsets of
\(\{2,3\}\).  Under \(2^x,3^x\in\Q\), every \(T_j^x\) is rational.
Therefore no condition asserting rationality or integrality of such
semisimple powers is stronger than the original hypotheses.
\end{proposition}

This rules out a broad class of seemingly richer finite-dimensional
constructions.  Any matrix-functional-calculus approach must introduce
a repeated node, a new algebraic spectral point, or a nonfunctional
relation such as addition.

\section{The fractional-part system and rational Koopman cocycles}

Return to \(x=n+\alpha\), \(0\leq\alpha<1\), and the rationals
\(r=2^\alpha\), \(s=3^\alpha\) from \eqref{eq:r-s}.  If \(\alpha=0\)
the conjecture is proved, so assume for analysis that \(\alpha\) is
irrational.

Let \(X=[0,1)\) with Lebesgue measure, and let
\[
        R_\alpha(t)=\{t+\alpha\}.
\]
Define the jump function
\[
        w(t)=\floor{t+\alpha}
        =\mathbf 1_{[1-\alpha,1)}(t).
\]
The elementary identity
\begin{equation}\label{eq:jump-coboundary}
        w(t)-\alpha=t-R_\alpha(t)
\end{equation}
is the source of both the usefulness and the limitations of the model.

For \(b=2,3\), define
\begin{equation}\label{eq:cocycle}
        c_b(t)=b^{w(t)-\alpha}.
\end{equation}
It takes exactly two values:
\[
        c_2(t)\in\left\{\frac1r,\frac2r\right\}\subset\Q,
        \qquad
        c_3(t)\in\left\{\frac1s,\frac3s\right\}\subset\Q.
\]
Conversely, rationality of either pair of values is equivalent to
rationality of \(b^\alpha\).

Let \(U_\alpha\) be the Koopman unitary
\[
        (U_\alpha\xi)(t)=\xi(R_\alpha(t))
\]
on \(L^2(X)\), and let
\[
        W_b=M_{c_b}U_\alpha.
\]

\begin{theorem}[The weighted operators are bounded coboundary twists]
\label{thm:koopman-similar}
Let \(h_b(t)=b^t\).  Then \(h_b\) and \(h_b^{-1}\) are bounded, and
\begin{equation}\label{eq:similarity}
        W_b=M_{h_b}U_\alpha M_{h_b}^{-1}.
\end{equation}
Consequently \(W_b\) is similar to \(U_\alpha\), and all spectral
invariants preserved by bounded similarity are independent of the
rationality of \(b^\alpha\).
\end{theorem}

\begin{proof}
By \eqref{eq:jump-coboundary},
\[
        c_b(t)=b^{t-R_\alpha(t)}
              =\frac{h_b(t)}{h_b(R_\alpha(t))}.
\]
Substitution into the definition of \(W_b\) gives
\eqref{eq:similarity}.
\end{proof}

For irrational \(\alpha\), \(U_\alpha\) has eigenvalues
\(\ee^{2\pi\ii k\alpha}\), \(k\in\Z\), dense in \(\mathbb T\), so
\[
        \Spec(W_b)=\Spec(U_\alpha)=\mathbb T.
\]
The iterated products telescope:
\begin{equation}\label{eq:cocycle-telescope}
 \prod_{j=0}^{N-1}c_b(R_\alpha^jt)
 =b^{t-R_\alpha^Nt},
\end{equation}
which is uniformly bounded above and below independently of \(N\).
Hence every logarithmic growth exponent is zero.  Also
\[
        \int_X\log c_b\,dt
        =(\log b)\int_X(w-\alpha)\,dt=0,
\]
so the Fuglede--Kadison determinant of the positive multiplier
\(M_{c_b}\) is \(1\).  The rotation has Kolmogorov--Sinai entropy zero
for every \(\alpha\) \cite{Walters}.  None of these invariants can detect whether
\(\alpha\) is rational.

\subsection{The common projection is the surviving datum}

Let
\[
        P=M_w.
\]
Then \(P\) is a projection and
\begin{equation}\label{eq:trace-alpha}
        \tau(P)=\int_Xw(t)\,dt=\alpha.
\end{equation}
The two rational multipliers are affine functions of the same
projection:
\begin{equation}\label{eq:weights-projection}
        M_{c_2}=\frac1r(I+P),\qquad
        M_{c_3}=\frac1s(I+2P).
\end{equation}
Equivalently,
\[
        2rM_{c_2}-sM_{c_3}=I.
\]
Thus the scalar hypotheses produce a rationally defined two-valued
observable, but its invariant trace may be irrational.  Rational
operator coefficients do not imply rational values of a trace state.

\begin{theorem}[Exact finite-dimensional trace criterion]
\label{thm:finite-trace}
Assume \eqref{eq:r-s}.  If there exist \(N\geq1\), a projection
\(P_N\in M_N(\C)\), and a trace-preserving identification of the
distinguished projection in the Koopman model such that
\[
        \frac1N\Tr(P_N)=\tau(P)=\alpha,
\]
then \(\alpha=0\), and hence \(x\in\Z\).
\end{theorem}

\begin{proof}
The normalized trace of a matrix projection is
\(\rank(P_N)/N\in\Q\), so \(\alpha\in\Q\).  By
\cref{lem:rational-collapse} and \(2^\alpha\in\Q\),
\(\alpha=0\).
\end{proof}

\begin{warning}[Approximation is not enough]
For every real \(\alpha\in[0,1]\), one can choose projections
\(P_N\in M_N(\C)\) of rank \(\floor{\alpha N}\), giving
\[
        \frac1N\Tr(P_N)\longrightarrow\alpha.
\]
Thus quasidiagonal, sofic, MF, or moment approximation that reproduces
the trace only in the limit cannot force rationality.  An operator proof
along this route needs an \emph{exact} finite-dimensional trace
realization, a bounded-denominator theorem, or an arithmetic restriction
on the admissible matrix sizes.
\end{warning}

\section{Finite Hankel realization of the Sturmian jump word}

Follow the orbit of \(0\) and define
\begin{equation}\label{eq:sturmian-word}
        w_k=w(R_\alpha^k0)
            =\floor{(k+1)\alpha}-\floor{k\alpha}
            \in\{0,1\}.
\end{equation}
This is the characteristic Sturmian word of slope \(\alpha\) when
\(\alpha\) is irrational \cite{Lothaire,Queffelec}.  It is recoverable from either rational
cocycle in \eqref{eq:cocycle}: the two values distinguish \(w_k=0\)
from \(w_k=1\).

Let
\[
        \mathsf H_w=(w_{i+j})_{i,j\geq0}
\]
be its infinite Hankel matrix.

\begin{theorem}[Finite Hankel rank closes the conjecture]
\label{thm:hankel-closes}
Under the Alaoglu--Erd\H{o}s hypotheses, if
\(\rank\mathsf H_w<\infty\) over \(\C\) (equivalently over any
characteristic-zero field containing the sequence), then
\(\alpha=0\), hence \(x\in\Z\).
\end{theorem}

\begin{proof}
Finite Hankel rank implies that \((w_k)\) satisfies a linear recurrence
of some finite order \cite{Kailath,Peller} \(d\), from some point onward.  Once the recurrence
is in force, the state
\[
        (w_k,w_{k+1},\ldots,w_{k+d-1})
\]
belongs to the finite set \(\{0,1\}^d\) and determines the next term.
Some state repeats, so the sequence is eventually periodic.

On the other hand, telescoping \eqref{eq:sturmian-word} gives
\[
 \frac1N\sum_{k=0}^{N-1}w_k
 =\frac{\floor{N\alpha}}{N}
 \longrightarrow\alpha.
\]
The limiting mean of an eventually periodic binary sequence is rational.
Thus \(\alpha\in\Q\), and \cref{lem:rational-collapse} gives
\(\alpha=0\).
\end{proof}

\begin{corollary}[No exact finite-dimensional linear realization]
\label{cor:no-finite-realization}
If \(x\) were a counterexample, there could be no finite-dimensional
linear system
\[
        w_k=\lambda(T^kv)
\]
with a fixed finite-dimensional vector space, operator \(T\), vector
\(v\), and linear functional \(\lambda\).  Any such realization would
give \(\rank\mathsf H_w<\infty\).
\end{corollary}

This criterion is more flexible than a finite-state automaton.  It rules
out every exact finite-dimensional linear recurrence realization,
including those over \(\R\) or \(\C\).  It suggests a concrete operator
target: derive finite Hankel rank, rationality of the generating function
\(\sum_{k\geq0}w_kz^k\), or a uniform bounded-denominator recurrence from
the simultaneous rationality of the two cocycles.  Ordinary low-rank
approximation is again insufficient: long finite Sturmian prefixes can
be approximated by periodic words of growing periods.


\section{Further operator routes: exact content and failure modes}

This section records several natural operator-theoretic ideas that were
tested against the local-to-global framework.  The purpose is to prevent
a reformulation from being mistaken for progress and to identify the
small modification that would turn each route into a genuine closing
argument.

\subsection{Wold theory and unilateral multiplication shifts}

On \(\ell^2(\N)\), the multiplication isometry
\(V_m e_n=e_{mn}\) is a pure isometry with countably infinite defect
multiplicity for every \(m\geq2\).  Hence individual operators \(V_m\)
are all unitarily equivalent by Wold theory.  A pair carries more orbit
information, but its bilateral dilation returns to the universality of
\cref{thm:bilateral-universal}.  Therefore neither the spectrum of one
\(V_m\) nor its Wold multiplicity can compare the labels \(m=2\) and
\(m=A\).

\emph{What would help:} retain the additive shift \(S\), the Hamiltonian
\(H\), or the concrete arithmetic basis.  The affine relation with \(S\)
is rigid by \cref{thm:affine-rigidity}; global \(H\)-covariance is rigid
by \cref{thm:global-isometry}.

\subsection{Tensor products and multiplicative closure}

Tensor powers of semisimple scalar modes produce eigenvalues
\(2^a3^b\).  Unique factorization makes all such labels distinct, so
the generated finite-dimensional functional calculus remains
semisimple.  Under the target map they become \(A^aB^b\), which are
again distinct.  Therefore tensoring cannot spontaneously create the
Jordan self-extension needed in \cref{thm:jordan-closes}.

\emph{What would help:} one additive decomposition that identifies two
different tensor constructions, one nontrivial extension class, or one
new simple object whose dimension is multiplicatively independent of
\(2,3\).

\subsection{Cuntz algebras and \(K\)-theory}

The Cuntz algebra \(\mathcal O_n\) has \cite{Cuntz1977}
\(K_0(\mathcal O_n)\cong\Z/(n-1)\Z\) and \(K_1(\mathcal O_n)=0\).
Thus a stable isomorphism or Morita equivalence between
\(\mathcal O_2\) and \(\mathcal O_A\) would force \(A=2\).  But the
scalar equality \(A=2^x\) supplies no such equivalence; postulating it
already inserts a strong bridge between the source and target labels.
At the other extreme, the \(K\)-theory of the fixed two-coordinate
Toeplitz algebra does not see the weights \(\log2,\log3\) at all.

\emph{What would help:} a naturally constructed correspondence or
bimodule whose \(K\)-theoretic class is forced by the scalar
exponentiation and simultaneously relates \(\mathcal O_2\) to
\(\mathcal O_A\).  Without that construction, the \(K\)-groups merely
restate that different known integers give different Cuntz algebras.

\subsection{Operator monotonicity and Loewner positivity}

For \(0<\alpha<1\), the function \(t^\alpha\) is operator monotone and
operator concave on \((0,\infty)\).  Its two-point Loewner matrix at
\(2,3\) is positive:
\[
 \begin{pmatrix}
   \alpha2^{\alpha-1}&3^\alpha-2^\alpha\\
   3^\alpha-2^\alpha&\alpha3^{\alpha-1}
 \end{pmatrix}\succeq0.
\]
Under \eqref{eq:r-s}, the off-diagonal entry and the factors
\(2^{\alpha-1},3^{\alpha-1}\) are rational; the unknown \(\alpha\)
remains in the diagonal derivative.  Positivity gives real inequalities,
not an arithmetic discreteness condition.

\emph{What would help:} algebraicity or rationality of any nonzero
diagonal derivative, the determinant, or a finite-rank degeneracy.
Each would make \(\alpha\) algebraic and then force \(\alpha=0\), but
none follows from positivity alone.

\subsection{Spectral determinants and heat asymptotics}

The local partition function \eqref{eq:partition}, its meromorphic
continuation, and the Weyl law \eqref{eq:weyl} depend analytically on
the continuous parameters \(\log a,\log b\).  Under
\((a,b)\mapsto(a^x,b^x)\) they transform by ordinary scaling.
Similarly, the Koopman cocycles have zero mean logarithm and bounded
products by \eqref{eq:cocycle-telescope}.  Determinants and Lyapunov
exponents therefore vanish or rescale exactly as they should for every
real \(x>0\).

\emph{What would help:} an external integral normalization of energy,
such as the full prime spectrum, together with an exact finite-lattice
implementation.  Without the integral lattice, asymptotic coefficients
are real-valued and cannot quantize the scale.

\subsection{Continuous interpolation of degree}

Maps \(z\mapsto z^n\) on the circle, endomorphisms of \(C(\mathbb T)\), and
finite-degree covering operators carry an integer topological degree.
A norm-continuous family cannot change this degree.  Consequently a
continuous ``fractional-degree'' interpolation from \(2\) to \(A\)
cannot remain inside the category of finite coverings.

This observation does not prove the conjecture: the hypotheses do not
provide such an interpolation.  It does show that attempts to realize
\(n\mapsto n^x\) as a continuous flow of genuine covering maps must
leave the arithmetic category between integer times.  The obstruction
is topological rather than transcendence-theoretic.

\subsection{Abstract dilation theory}

The two-mode embedding \(J\) in \eqref{eq:fock-embedding} is an
isometry, and abstract isometries admit many unitary or isometric
dilations.  Such dilations can assign arbitrary vectors to the
orthogonal complement and therefore erase the valuation lattice.
They do not imply that \(p^x\) is rational for a new prime \(p\).

\emph{What would help:} a dilation theorem inside a rigid category of
monomial operators, product systems with integral degree maps, graph
correspondences with finite incidence matrices, or based Hilbert
modules whose basis is arithmetically distinguished.  In any finite
such category, \cref{thm:finite-prime-closure} becomes the closing
mechanism.

\section{Five exact closing bridges}

The operator analysis reduces a large search space to five compact
targets.  Each target adds information of a different kind, but each is
provably sufficient.

\subsection{Bridge I: one third algebraic mode}

By \cref{thm:three-mode}, prove \(q^x\in\Alg\) for one
\(q\in\Qp\setminus\langle2,3\rangle\).  In Hamiltonian language, extend
the partial intertwiner to one additional eigenvector:
\[
        Te_q\neq0,\qquad HT=xTH,
\]
with \(Te_q\) supported at an algebraic arithmetic label.  This is the
smallest propagation target and uses the full strength of the Six
Exponentials Theorem immediately.

\subsection{Bridge II: one rational first jet}

Construct any rational square-zero perturbation \(N\neq0\) for which
\((2I+N)^x\) or \((3I+N)^x\) is rational.  By
\cref{thm:jordan-closes}, this makes \(x\) rational.  Conceptually, this
requires a derivation or self-extension compatible with the scalar
power operation.

A potentially useful algebraic formulation is to seek a
\(\Q\)-derivation \(\delta\) on an operator algebra containing the
scalar mode such that
\[
        \delta(a^x)=x a^{x-1}\delta(a)
\]
holds and both \(\delta(a)\) and \(\delta(a^x)\) are rationally
controlled.  If \(\delta(2)\neq0\), the chain rule gives the missing
coefficient.

\subsection{Bridge III: finite prime-mode completion}

Extend the two valuation columns
\[
        \nu(2)\longmapsto\nu(A),\qquad
        \nu(3)\longmapsto\nu(B)
\]
to an integral endomorphism of a finite prime lattice that scales the
log-energy vector.  Then \cref{thm:finite-prime-closure} applies.  A
unimodular extension is even more rigid and forces \(x=1\).

This target is natural for finite graph correspondences, substitution
matrices, fusion rings, and finite product systems.  The audit question
for any proposed construction is simple: does its incidence matrix
really satisfy \(M^{\mathsf T}\ell=x\ell\), or does it merely reproduce
the two known columns?

\subsection{Bridge IV: exact finite realization of the jump observable}

Produce an exact finite-dimensional tracial model for \(P\) or an exact
finite-dimensional linear realization of the Sturmian sequence
\((w_k)\).  The former invokes \cref{thm:finite-trace}; the latter invokes
\cref{thm:hankel-closes}.  Approximation with growing dimensions or
periods does not suffice.

A more operator-invariant variant uses the correlation sequence
\[
        c_k=\tau(PU_\alpha^kPU_\alpha^{-k}).
\]
Its Fourier representation contains infinitely many nonzero frequencies
when \(\alpha\) is irrational.  Any theorem forcing a finite spectral
support, rational generating function, or bounded-denominator exact
recurrence for \((c_k)\) would likewise force rationality of
\(\alpha\).

\subsection{Bridge V: one additive--multiplicative coherence}

Construct a morphism preserving one affine relation such as
\(V_2S=S^2V_2\), or one categorical relation such as
\(E_2\otimes E_2\cong E_2\oplus E_2\).  The robust gap in
\cref{thm:affine-rigidity} shows that even preservation modulo compact
operators is enough in the concrete affine model.

A weaker propagation version would only need to manufacture a third
mode.  For example, a rule that makes the label \(5=2+3\) available
after exponentiation would give algebraicity of \(5^x\), and
\cref{thm:three-mode} would finish the argument before full additive
functoriality is required.

\section{Prioritized research program}

The following program is ordered by the amount of genuinely new
arithmetic structure required, not by aesthetic appeal.

\subsection{Priority 1: third-mode leakage from an additive transfer operator}

The largest payoff per lemma is to derive algebraicity of \(5^x\) (or
any non-\(3\)-smooth rational base).  The affine semigroup suggests
studying transfer operators that mix translation and dilation:
\[
        (L_mf)(k)=\sum_{j=0}^{m-1}f\!\left(\frac{k+j}{m}\right)
\]
in a suitable discrete or compact model, together with the
multiplication isometries \(V_2,V_3\).  The goal is not a full
endomorphism of the \(ax+b\) algebra---which would already force
\(x=1\)---but a weaker positivity or integrality statement that produces
one new algebraic mode.

A concrete test is whether the ranges or defect projections attached to
\(V_2\) and \(V_3\) can be combined canonically to create the degree-five
range projection.  Any construction must be audited against
\cref{thm:bilateral-universal}: if it uses only the commuting pair and
not the additive placement of their ranges, it cannot distinguish the
labels.

\subsection{Priority 2: rational self-extension from a deformation}

Seek a rational one-parameter deformation \(T(t)\) through \(T(0)=2I\)
such that \(T(t)^x\) is rational to first order at \(t=0\).  Differentiating
would give
\[
        \frac{d}{dt}T(t)^x\bigg|_{t=0}
        =x2^{x-1}T'(0),
\]
and any nonzero rational \(T'(0)\) closes the problem.

Possible sources include:
\begin{itemize}
\item a derivation on a rational semigroup algebra;
\item a nontrivial extension in a category generated by the two scalar
      modes and an arithmetic correspondence;
\item a deformation of a transfer operator whose matrix entries remain
      rational because they count finite fibers;
\item a divided-difference construction in which two spectral nodes are
      forced to coalesce without leaving the rational category.
\end{itemize}
The semisimple theorem shows exactly what must be avoided: coupling only
the distinct nodes \(2\) and \(3\) gives \(B-A\), never a derivative.

\subsection{Priority 3: finite correspondence or fusion completion}

Build a finite based module with basis indexed by a finite prime set and
an integral endomorphism extending the two known valuation columns.
The logarithmic dimension vector should be a left eigenvector with
eigenvalue \(x\).  Once this is achieved, no delicate spectral estimate
is needed; \cref{thm:finite-prime-closure} and Gelfond--Schneider finish
the proof.

The hard point is not finding an integer matrix with the two desired
columns; arbitrary completion is easy.  The hard point is forcing all
columns to obey the same logarithmic scaling.  Candidate mechanisms are
associativity constraints, duals in a rigid tensor category, finite
index, or a trace/dimension function compatible with both source and
target.

\subsection{Priority 4: exact Hankel rigidity from two rational cocycles}

The pair of rational-valued cocycles is stronger than one such cocycle
only through the common exponent:
\[
        c_2=2^{P-\alpha I},\qquad
        c_3=3^{P-\alpha I}.
\]
Both are affine rational functions of the same projection \(P\), so
purely algebraic elimination recovers \(P\) but not its trace.  A
successful argument must constrain the orbit of \(P\), not merely \(P\)
itself.

Promising exact statements would be:
\begin{enumerate}
\item a bounded-order recurrence for the words obtained from both
      cocycles;
\item a uniform denominator bound for all correlation moments;
\item an exact rational generating function for one nontrivial orbit
      observable;
\item a finite-dimensional representation whose dimension is bounded
      independently of the approximation scale.
\end{enumerate}
Any of these would invoke \cref{thm:hankel-closes} or
\cref{thm:finite-trace}.

\subsection{Priority 5: quantitative affine defects}

The gap
\(\|V_AS^d-S^{2d}V_A\|_{\ess}\geq\sqrt2\) for \(A\neq2\) suggests
looking for an arithmetic construction whose affine defect can be shown
strictly smaller than \(\sqrt2\).  This is weaker than exact
functoriality and robust under compact errors.  Similar estimates can be
studied in normalized Schatten norms on finite truncations.  In the
standard basis, the wrong label produces a defect of fixed density, not
a boundary effect.

The challenge is to connect the scalar equalities \(A=2^x,B=3^x\) to
such a norm estimate.  Heat-kernel averaging or transfer-operator
regularization might turn logarithmic scaling into approximate affine
coherence, but any estimate must beat a fixed, representation-independent
gap rather than merely tend to a nonzero constant.

\section{Logical audit: what has and has not been proved}

The conjecture itself remains open in this report.  The rigorous gains
are structural and conditional:

\begin{enumerate}
\item The largest rational/algebraic sector available under a
      counterexample is exactly the rank-two group generated by \(2\)
      and \(3\) (\cref{thm:three-mode}).
\item Pure bilateral shift data is completely universal
      (\cref{thm:bilateral-universal}).
\item Local Hamiltonian, heat, and KMS data is exactly related by time
      rescaling (\cref{prop:fock-homothety}).
\item A global monomial Hamiltonian extension forces a positive integer
      exponent (\cref{thm:global-isometry}).
\item A finite invariant prime lattice forces an integer exponent, and
      a lattice automorphism forces \(x=1\)
      (\cref{thm:finite-prime-closure,cor:automorphic-closure}).
\item One affine relation is robustly label-rigid, even modulo compact
      operators (\cref{thm:affine-rigidity}).
\item Semisimple rational functional calculus at \(2,3\) is automatic
      and cannot help, while one Jordan self-extension forces
      rationality (\cref{thm:semisimple-blind,thm:jordan-closes}).
\item The Koopman cocycles are bounded coboundaries; exact finite trace
      or Hankel realization closes, but approximation does not
      (\cref{thm:koopman-similar,thm:finite-trace,thm:hankel-closes}).
\end{enumerate}

None of the additional structures in the right-hand column below has
been derived from the scalar hypotheses.

\begin{center}
\begin{tabular}{p{0.42\textwidth}p{0.47\textwidth}}
\toprule
\textbf{Given unconditionally} & \textbf{Still missing}\\
\midrule
Two integer labels \(A=2^x,B=3^x\) &
One third independent algebraic label \(q^x\).\\
A rank-two covariant Fock embedding &
A finite or global integral lattice completion.\\
All semisimple rational two-point lifts &
One rational nonsemisimple self-extension.\\
Two rational-valued bounded coboundaries &
An exact finite trace or finite-Hankel realization.\\
Pure multiplicative semigroup coherence &
One additive--multiplicative coherence relation.\\
\bottomrule
\end{tabular}
\end{center}

This distinction is essential.  Each missing item is a genuine
strengthening, not a disguised consequence already established in the
report.

\section{Updated literature reconnaissance}

The literature search for this continuation was organized around three
separate questions:

\begin{enumerate}
\item Has the special rational/integer case of the Four Exponentials
      problem, or the Alaoglu--Erd\H{o}s conjecture itself, been resolved?
\item Have recent algebraic-independence results lowered the unconditional
      \(2\times3\) Six Exponentials threshold to the needed \(2\times2\)
      array in this arithmetic setting?
\item Do recent semigroup \(C^*\)-algebra, arithmetic-dynamics, or
      operator-algebra constructions supply a canonical additive or
      finite-lattice extension of the two multiplicative modes?
\end{enumerate}

Searches were run through 21 August 2026 using the exact phrases
``Alaoglu--Erdos conjecture'', ``four exponentials conjecture'',
``six exponentials theorem'', ``Schanuel conjecture'', ``algebraic
independence logarithms'', ``Bost--Connes'', ``ax+b semigroup
\(C^*\)-algebra'', and combinations with ``arXiv''.  Primary sources,
arXiv records, journal pages, and author-hosted surveys were preferred.

The search found no verified unconditional resolution of the conjecture
or of the Four Exponentials Conjecture.  The operative unconditional
threshold remains the Six Exponentials Theorem used in
\cref{thm:three-mode}.  Modern work on Schanuel-type statements and
exponential algebraic closedness \cite{Kirby2010,WaldschmidtOpen},, period conjectures, explicit linear
forms in logarithms, and zero estimates refines surrounding theory, but
does not turn the particular vanishing determinant
\[
        (\log2)(\log B)-(\log3)(\log A)=0
\]
into a contradiction.  Linear-form estimates apply when a nonzero
linear form in logarithms must be bounded away from zero; here the hard
relation is bilinear in logarithms and is exactly zero.

On the operator side, the Bost--Connes system, semigroup
\(C^*\)-algebras of \(ax+b\) type, product systems, and boundary
quotients provide a mature language for the logarithmic Hamiltonian and
for additive--multiplicative coherence.  They support the distinction
proved here: multiplicative time evolution is naturally homogeneous,
whereas additive range relations remember integer degree.  No source
located in the search derives the required additive or finite-prime
extension from the two scalar algebraicity assumptions alone.

The reproducible arXiv metadata ledger below records recent items
returned by the focused queries.  It is included to make the search
auditable rather than to imply that every listed paper directly attacks
the conjecture.

\begin{longtable}{>{\raggedright\arraybackslash}p{0.10\textwidth}
                  >{\raggedright\arraybackslash}p{0.23\textwidth}
                  >{\raggedright\arraybackslash}p{0.48\textwidth}
                  >{\raggedright\arraybackslash}p{0.10\textwidth}}
\toprule
\textbf{Date} & \textbf{Authors} & \textbf{Title} & \textbf{arXiv}\\
\midrule
\endfirsthead
\toprule
\textbf{Date} & \textbf{Authors} & \textbf{Title} & \textbf{arXiv}\\
\midrule
\endhead
\multicolumn{4}{p{0.94\textwidth}}{\emph{The live arXiv Atom harvest was not available inside the compilation sandbox.  The search protocol and stable references in the bibliography remain reproducible; no unverified recent title has been inserted.}}\\
\bottomrule
\end{longtable}

\begin{remark}[How to interpret the ledger]
The phrase searches inevitably return papers about model theory,
exponential fields, linear forms, or operator systems that are adjacent
rather than directly applicable.  Their relevance is negative as well
as positive: they confirm that the known unconditional theorem supplies
three logarithmic columns, while the present conjecture supplies two.
No claim of completeness is made for indexing services that were not
openly searchable.
\end{remark}


\section{Conclusions}

Operator theory does not bypass the transcendence barrier merely by
placing the four algebraic numbers \(2,3,A,B\) on a Hilbert space.  In
the natural two-mode representations, it explains why the barrier is so
stable: all multiplicative shift pairs are universal after bilateral
amplification, and all local logarithmic Hamiltonians are related by a
continuous time homothety.  Spectra, heat traces, local equilibrium
states, entropy, and bounded-coboundary invariants are therefore exactly
as compatible with a nonintegral \(x\) as with an integral one.

The useful operator contribution is a precise map of where discreteness
re-enters.

\begin{itemize}
\item The full arithmetic Hamiltonian has simple, discretely labelled
      spectrum.  A global isometric rescaling must send every \(n\) to
      the integer \(n^x\), and finite differences force \(x\in\N\).
\item A finite prime-mode completion turns logarithmic scaling into an
      integer-matrix eigenvalue problem.  Gelfond--Schneider then forces
      the eigenvalue to be an integer.
\item The affine relation \(V_mS=S^mV_m\) remembers the integer label in
      a way the multiplicative pair cannot; even a compact perturbation
      cannot disguise a wrong label.
\item Semisimple functional calculus sees only the two given values.
      A Jordan self-extension sees a derivative, and that first jet
      immediately rationalizes \(x\).
\item The Koopman model converts the fractional part into the trace of
      one projection.  Exact finite-dimensional or finite-Hankel
      realization discretizes the trace; approximation does not.
\item Six Exponentials shows that a hypothetical counterexample has a
      maximal algebraic sector of rank exactly two.  One new independent
      algebraic mode is enough.
\end{itemize}

The most promising next attack is therefore not ``find a subtler
spectrum.''  It is to derive one small coherence bridge that cannot be
implemented in the universal rank-two model: a third mode, a first jet,
a finite integral completion, an exact finite realization, or an
additive range relation.  These targets are narrow enough to test
individual constructions decisively, yet strong enough that no further
transcendence breakthrough would be needed once any one of them is
obtained.

\appendix

\section{Domain details for arithmetic Hamiltonians}

This appendix records the unbounded-operator points used implicitly in
the main text.

Let \(H\) be diagonal on \(\ell^2(\N)\),
\(He_n=(\log n)e_n\), with maximal domain
\[
 \mathcal D(H)
 =\left\{\xi=(\xi_n):
   \sum_{n\geq1}(\log n)^2|\xi_n|^2<\infty\right\}.
\]
The finite-support subspace
\[
        \mathcal C=\spn\{e_n:n\in\N\}
\]
is a core for \(H\).  If a bounded operator \(T\) satisfies
\(HT=cTH\) on \(\mathcal C\), then \(T\mathcal C\subseteq\mathcal D(H)\).
For each \(n\), spectral simplicity gives
\[
        Te_n\in\ker(H-c\log n).
\]
Thus either \(Te_n=0\), or \(n^c\in\N\) and
\(Te_n=\lambda_ne_{n^c}\) for some scalar \(\lambda_n\).  In
particular:

\begin{proposition}\label{prop:all-H-intertwiners}
Every bounded \(H\)-covariant operator has the form
\[
        Te_n=
        \begin{cases}
        \lambda_ne_{n^c},&n^c\in\N,\\
        0,&n^c\notin\N,
        \end{cases}
\]
with a bounded scalar family \((\lambda_n)\).  It is an isometry exactly
when \(n^c\in\N\) and \(|\lambda_n|=1\) for every \(n\).
\end{proposition}

This proposition explains why support, not spectral multiplicity, is the
central invariant.  Under a hypothetical counterexample, the support of
every such algebraically labelled intertwiner is confined by
\cref{thm:three-mode} to the \(3\)-smooth sector.

On \(\ell^2(\Qp)\), the same argument applies to
\(H_Ge_q=(\log q)e_q\).  A covariant operator is monomial on the rational
basis wherever \(q^c\in\Qp\).  Under a nonintegral
Alaoglu--Erd\H{o}s solution, \cref{thm:three-mode} says the algebraic
support is precisely \(\langle2,3\rangle\).

\section{The finite-difference integrality lemma in greater generality}

The proof of \cref{thm:global-isometry} used a standard but useful
criterion.

\begin{theorem}[Integer-valued real powers]\label{thm:integer-valued-powers}
Let \(c\geq0\).  If \(n^c\in\Z\) for every sufficiently large integer
\(n\), then \(c\in\Nzero\).
\end{theorem}

\begin{proof}
If \(c\notin\Nzero\), put \(r=\floor c+1\).  For all sufficiently large
\(n\), every term \((n+j)^c\), \(0\leq j\leq r\), is an integer, hence
\(\Delta^r n^c\in\Z\).  The repeated-integral identity
\[
 \Delta^r f(n)
 =\int_{[0,1]^r}f^{(r)}(n+t_1+\cdots+t_r)
   \,dt_1\cdots dt_r
\]
holds for \(f\in C^r\).  With \(f(t)=t^c\), the derivative is
\[
        f^{(r)}(t)
        =c(c-1)\cdots(c-r+1)t^{c-r}.
\]
It is nonzero with fixed sign on \((0,\infty)\), while \(c-r<0\).
Therefore \(\Delta^r n^c\) is nonzero and tends to zero, contradicting
its eventual integrality.
\end{proof}

The theorem remains valid if \(n^c\) is required to lie in any discrete
subset of \(\R\) with a uniform separation near the relevant range,
provided the forward differences inherit that discreteness.  This is
one reason exact lattice-valued operator models are powerful, while
asymptotic closeness to a lattice is not.

\section{Finite prime lattices as based operator correspondences}

Let \(P=\{p_1,\ldots,p_d\}\) and let
\(\ell=(\log p_1,\ldots,\log p_d)^{\mathsf T}\).
An integer matrix \(M\) defines a monomial map on Laurent monomials:
\[
 z^v\longmapsto z^{Mv},\qquad v\in\Z^d.
\]
On the group Hilbert space \(\ell^2(\Z^d)\), it defines
\(T_Me_v=e_{Mv}\) when \(M\) is injective.  The energy relation is
\[
 H_\ell T_Me_v
 =\langle\ell,Mv\rangle e_{Mv}
 =\langle M^{\mathsf T}\ell,v\rangle e_{Mv}.
\]
Hence \(H_\ell T_M=cT_MH_\ell\) exactly when
\(M^{\mathsf T}\ell=c\ell\).

The same matrix can be viewed in several equivalent categories:

\begin{enumerate}
\item as an endomorphism of the free abelian multiplicative group
      \(\Gamma_P\);
\item as a monomial endomorphism of the algebraic torus
      \((\C^\times)^d\);
\item as a based correspondence on the valuation lattice;
\item when \(M\geq0\), as a multi-shift on the positive Fock cone;
\item as an incidence matrix with logarithmic dimension vector
      \(\ell\).
\end{enumerate}

In every finite realization, the scaling factor is an eigenvalue of a
monic integer polynomial.  This is the categorical content of
\cref{thm:finite-prime-closure}.  Passing to an infinite prime set
removes the algebraic-integer conclusion: an infinite integer matrix may
have arbitrary positive eigenvalues in an appropriate functional-analytic
sense.  Therefore finiteness is not cosmetic; it is the source of
arithmetic discreteness.

\subsection{A completion test}

Suppose a construction proposes a finite matrix \(M\) extending the two
known columns.  The following checklist detects circularity.

\begin{enumerate}
\item Are the columns genuinely integral valuation vectors?
\item Does \(M^{\mathsf T}\ell=x\ell\) hold for \emph{every} column, or
      only for the columns indexed by \(2\) and \(3\)?
\item Is the prime set finite and invariant under all introduced power
      images?
\item Was rationality of a new \(p^x\) proved, or silently assumed when
      its valuation column was written down?
\item If \(M\) is claimed unimodular, is surjectivity of the power map
      actually available?
\end{enumerate}

A positive answer to the first three questions, without assuming the
desired conclusion in the fourth, is already a proof.

\section{Functional calculus on primary components}

Let \(T\) be a rational matrix with positive spectrum contained in
\(\{2,3\}\).  Its minimal polynomial has the form
\[
        m_T(t)=(t-2)^r(t-3)^s.
\]
Hermite interpolation gives a polynomial \(P_x(t)\) of degree
less than \(r+s\) such that
\[
 P_x^{(j)}(2)=\frac{d^j}{dt^j}t^x\bigg|_{t=2},
 \quad 0\leq j<r,
\]
and similarly at \(3\).  Then \(T^x=P_x(T)\).  Consequently:

\begin{proposition}[Jet hierarchy]\label{prop:jet-hierarchy}
Under \(2^x,3^x\in\Q\), rationality of \(T^x\) is automatic for every
semisimple \(T\) (\(r,s\leq1\)).  If \(r\geq2\) or \(s\geq2\), it
additionally requires rational control of at least one derivative
\[
        (x)_j\,2^{x-j}
        \quad\text{or}\quad
        (x)_j\,3^{x-j},
\]
where \((x)_j=x(x-1)\cdots(x-j+1)\).
The first nonzero self-extension, \(j=1\), already forces \(x\in\Q\).
\end{proposition}

This hierarchy precisely classifies the information visible to every
finite-dimensional rational functional-calculus model supported at the
two original bases.

\subsection{Why the distinct-node triangular matrix is split}

For completeness, the matrix
\[
        T_c=\begin{pmatrix}2&c\\0&3\end{pmatrix}
\]
is rationally diagonalizable:
\[
 \begin{pmatrix}1&c\\0&1\end{pmatrix}^{-1}
 T_c
 \begin{pmatrix}1&c\\0&1\end{pmatrix}
 =
 \begin{pmatrix}2&0\\0&3\end{pmatrix}
\]
after the appropriate sign convention for the off-diagonal conjugating
entry.  More invariantly, the coprimality of \(t-2\) and \(t-3\) makes
the module split.  Its off-diagonal entry in \(T_c^x\) is a divided
difference, not a derivative.  A genuine Jordan block cannot be reached
by a rational similarity while keeping distinct eigenvalues.

\section{Koopman calculations and correlation data}

The jump identity
\[
        \sum_{j=0}^{N-1}(w(R_\alpha^jt)-\alpha)
        =t-R_\alpha^Nt
\]
proves bounded discrepancy directly.  In particular, the Sturmian word
has balanced prefixes:
\[
        \left|\sum_{j=0}^{N-1}w_j-N\alpha\right|<1.
\]
This strong regularity does not imply periodicity; it is shared by every
irrational rotation word.

The distinguished interval projection has Fourier coefficients
\[
 \widehat w(k)
 =
 \begin{cases}
 \alpha,&k=0,\\[0.6ex]
 \displaystyle
 \frac{\ee^{-2\pi\ii k(1-\alpha)}-\ee^{-2\pi\ii k}}{2\pi\ii k},
 &k\neq0.
 \end{cases}
\]
For irrational \(\alpha\), infinitely many coefficients are nonzero.
The correlation sequence
\[
 c_n=\tau(PU_\alpha^nPU_\alpha^{-n})
\]
has the Fourier expansion
\[
        c_n=\sum_{k\in\Z}|\widehat w(k)|^2
             \ee^{2\pi\ii kn\alpha}.
\]
Thus its spectral measure has infinitely many atoms for irrational
\(\alpha\).  Any exact representation that forces only finitely many
frequencies would make \(\alpha\) rational.  This is the correlation
counterpart of the Hankel-rank criterion.

\subsection{Exact versus asymptotic matrix models}

There are three materially different claims:

\begin{enumerate}
\item \emph{Exact fixed-size model:} a single \(M_N\) reproduces the
      projection and its trace exactly.  This forces
      \(\alpha\in\frac1N\Z\).
\item \emph{Uniformly bounded size:} a sequence of exact models has
      dimensions bounded by \(N_0\).  Passing to a constant-size
      subsequence again forces a rational trace with bounded
      denominator.
\item \emph{Growing-size approximation:} dimensions tend to infinity
      and traces converge to \(\alpha\).  This is compatible with every
      real \(\alpha\) and gives no arithmetic conclusion.
\end{enumerate}

Any future claim based on microstates or finite sections should specify
which of these regimes is proved.

\section{Dependency map for the closing criteria}

The logical flow can be summarized without operator terminology.

\begin{align*}
\text{third independent algebraic mode}
&\Longrightarrow x\in\Q
 &&\text{(Six Exponentials)},\\
\text{rational first derivative at \(2\) or \(3\)}
&\Longrightarrow x\in\Q
 &&\text{(elementary division)},\\
\text{finite invariant prime lattice}
&\Longrightarrow x\text{ algebraic integer}
 &&\text{(integer matrix)},\\
x\text{ algebraic and }2^x\in\Alg
&\Longrightarrow x\in\Q
 &&\text{(Gelfond--Schneider)},\\
x\in\Q,\ 2^x\in\Q
&\Longrightarrow x\in\Z
 &&\text{(valuation)},\\
\text{exact finite trace/Hankel model}
&\Longrightarrow\alpha\in\Q
 &&\text{(finite rank/periodicity)},\\
\alpha\in\Q,\ 2^\alpha\in\Q
&\Longrightarrow\alpha=0
 &&\text{(valuation)},\\
\text{one affine coherence}
&\Longrightarrow A=2
 &&\text{(translation length)},\\
A=2^x=2
&\Longrightarrow x=1.
\end{align*}

This map makes the role of each external theorem explicit and prevents
a closing criterion from being confused with a derivation of its
premise.

\section{Reproducible search notes}

The arXiv metadata queries used for the ledger were variants of:

\begin{verbatim}
all:"four exponentials conjecture"
(all:Schanuel OR all:"six exponentials" OR
 all:"four exponentials") AND cat:math.NT
(all:"algebraic independence" AND all:logarithms) AND cat:math.NT
(all:"Bost-Connes" OR all:"ax+b semigroup" OR
 all:"semigroup C*-algebra") AND cat:math.OA
\end{verbatim}

Results were sorted by submission date and screened for direct
relevance.  General web searches additionally used the exact conjecture
name with the scalar condition \(2^x,3^x\in\Z\), and checked cited
surveys and journal versions.  Because negative literature searches can
never prove nonexistence, the report's mathematical claims rely only on
the theorems proved here or on standard cited results, not on the
absence of a search hit.

\begin{thebibliography}{99}

\bibitem{AlaogluErdos1944}
L.~Alaoglu and P.~Erd\H{o}s,
\newblock On highly composite and similar numbers,
\newblock \emph{Transactions of the American Mathematical Society}
  \textbf{56} (1944), 448--469.

\bibitem{BakerBook}
A.~Baker,
\newblock \emph{Transcendental Number Theory},
\newblock second edition, Cambridge University Press, 1990.

\bibitem{BostConnes}
J.-B.~Bost and A.~Connes,
\newblock Hecke algebras, type III factors and phase transitions with
spontaneous symmetry breaking in number theory,
\newblock \emph{Selecta Mathematica (N.S.)} \textbf{1} (1995), 411--457.

\bibitem{BratteliRobinson}
O.~Bratteli and D.~W.~Robinson,
\newblock \emph{Operator Algebras and Quantum Statistical Mechanics II},
\newblock second edition, Springer, 1997.

\bibitem{BrownOzawa}
N.~P.~Brown and N.~Ozawa,
\newblock \emph{\(C^*\)-Algebras and Finite-Dimensional Approximations},
\newblock Graduate Studies in Mathematics 88, American Mathematical
Society, 2008.

\bibitem{BrownloweEtAl}
N.~Brownlowe, A.~an Huef, M.~Laca, and I.~Raeburn,
\newblock Boundary quotients of the Toeplitz algebra of the affine
semigroup over the natural numbers,
\newblock \emph{Ergodic Theory and Dynamical Systems} \textbf{32} (2012),
35--62.

\bibitem{Cuntz1977}
J.~Cuntz,
\newblock Simple \(C^*\)-algebras generated by isometries,
\newblock \emph{Communications in Mathematical Physics} \textbf{57}
(1977), 173--185.

\bibitem{CuntzAxB}
J.~Cuntz,
\newblock \(C^*\)-algebras associated with the \(ax+b\)-semigroup over
\(\N\),
\newblock in \emph{K-Theory and Noncommutative Geometry}, EMS Series of
Congress Reports, European Mathematical Society, 2008, 201--215;
\newblock \url{https://arxiv.org/abs/math/0611541}.

\bibitem{CuntzDeningerLaca}
J.~Cuntz, C.~Deninger, and M.~Laca,
\newblock \(C^*\)-algebras of Toeplitz type associated with algebraic
number fields,
\newblock \emph{Mathematische Annalen} \textbf{355} (2013), 1383--1423.

\bibitem{Gelfond}
A.~O.~Gelfond,
\newblock \emph{Transcendental and Algebraic Numbers},
\newblock Dover Publications, 1960.

\bibitem{Higham}
N.~J.~Higham,
\newblock \emph{Functions of Matrices: Theory and Computation},
\newblock Society for Industrial and Applied Mathematics, 2008.

\bibitem{Kailath}
T.~Kailath,
\newblock \emph{Linear Systems},
\newblock Prentice--Hall, 1980.

\bibitem{Kirby2010}
J.~Kirby,
\newblock Schanuel's conjecture and the exponential algebraic closedness
of the complex field,
\newblock \emph{Bulletin of the London Mathematical Society}
\textbf{42} (2010), 222--234;
\newblock \url{https://arxiv.org/abs/0810.4285}.

\bibitem{LacaRaeburn}
M.~Laca and I.~Raeburn,
\newblock Semigroup crossed products and the Toeplitz algebras of
nonabelian groups,
\newblock \emph{Journal of Functional Analysis} \textbf{139} (1996),
415--440.

\bibitem{Lang}
S.~Lang,
\newblock \emph{Introduction to Transcendental Numbers},
\newblock Addison--Wesley, 1966.

\bibitem{LindMarcus}
D.~Lind and B.~Marcus,
\newblock \emph{An Introduction to Symbolic Dynamics and Coding},
\newblock Cambridge University Press, 1995.

\bibitem{Lothaire}
M.~Lothaire,
\newblock \emph{Algebraic Combinatorics on Words},
\newblock Encyclopedia of Mathematics and its Applications 90,
Cambridge University Press, 2002.

\bibitem{Nica}
A.~Nica,
\newblock \(C^*\)-algebras generated by isometries and Wiener--Hopf
operators,
\newblock \emph{Journal of Operator Theory} \textbf{27} (1992), 17--52.

\bibitem{Peller}
V.~V.~Peller,
\newblock \emph{Hankel Operators and Their Applications},
\newblock Springer Monographs in Mathematics, Springer, 2003.

\bibitem{Queffelec}
M.~Queff\'elec,
\newblock \emph{Substitution Dynamical Systems---Spectral Analysis},
\newblock second edition, Lecture Notes in Mathematics 1294, Springer,
2010.

\bibitem{Schneider}
T.~Schneider,
\newblock \emph{Einf\"uhrung in die transzendenten Zahlen},
\newblock Springer, 1957.

\bibitem{WaldschmidtBook}
M.~Waldschmidt,
\newblock \emph{Diophantine Approximation on Linear Algebraic Groups:
Transcendence Properties of the Exponential Function in Several
Variables},
\newblock Grundlehren der mathematischen Wissenschaften 326, Springer,
2000.

\bibitem{WaldschmidtOpen}
M.~Waldschmidt,
\newblock Open Diophantine problems,
\newblock \emph{Moscow Mathematical Journal} \textbf{4} (2004), 245--305;
\newblock \url{https://arxiv.org/abs/math/0312440}.

\bibitem{Walters}
P.~Walters,
\newblock \emph{An Introduction to Ergodic Theory},
\newblock Graduate Texts in Mathematics 79, Springer, 1982.

\end{thebibliography}

\end{document}
